% !TEX root = ../main.tex
\chapter{Direct and Indirect Finance}\label{ch:directindirect}
\begin{paracol}{2}

\begin{sources}
\begin{tabularx}{\linewidth}{@{}l l L@{}}
\toprule
\textbf{Segment} & \textbf{Length} & \textbf{Content}\\
\midrule
\seg{17}{1}{L17.1 FT: Shadow banking} & 4:15 & Adair Turner: shadow banking is like cholesterol; traditional versus shadow bank.\\
\seg{17}{2}{L17.2 Bagehot's World} & 8:51 & Money market and capital market as separate worlds.\\
\seg{17}{3}{L17.3 The New World} & 10:55 & American banks: long-term assets, interbank repo, the origin of rating agencies.\\
\seg{17}{4}{L17.4 Funding liquidity versus market liquidity} & 2:33 & Moulton's shiftability (1918).\\
\seg{17}{5}{L17.5 Digression: Schumpeter on banking} & 4:05 & Banks as engines of development.\\
\seg{17}{6}{L17.6 Payment versus funding} & 5:31 & A loan funded by deposits, then taken out by a bond.\\
\seg{17}{7}{L17.7 Reading: Gurley and Shaw} & 2:04 & \emph{Money in a Theory of Finance} (1960): direct and indirect finance.\\
\seg{17}{8}{L17.8 Financial evolution} & 13:12 & Pension funds and insurers; from intermediation to mutual funds.\\
\seg{17}{9}{L17.9 Banking evolution} & 11:12 & Deposits as cheap permanent funding; from Jimmy Stewart banks to shadow banks.\\
\seg{17}{10}{L17.10 Preview: Central banking and shadow banking} & 8:28 & Liquidity puts, solvency versus liquidity, market liquidity provided by dealers.\\
\bottomrule
\end{tabularx}

\smallskip
\textbf{Files.} Transcripts \file{materials/transcripts/L17-P*.txt}; Mehrling's notes \mn{17}{1}.
\textbf{Reading.} J.~G.~Gurley and E.~S.~Shaw, \emph{Money in a Theory of Finance}, Brookings Institution, 1960 (passages).
\end{sources}

\section*{Motivation}
This lecture opens the fourth and last part of the course, which extends the money view from the money market to the capital market, to asset prices and the
prices of risk \ts{17}{1}{0:49}\ts{17}{2}{0:03}. The two subjects are usually taught separately: capital markets and CAPM in finance courses in the
business school, money and banking (perhaps IS--LM) in economics departments \ts{17}{2}{0:03}. Shadow banking shows that this
intellectual separation does not match the real world, where the two are more intertwined than at any time since World War~II, perhaps World War~I
\ts{17}{2}{1:04}. The separation has a historical origin, which the lecture reconstructs.

% ------------------------------------------------------------------
\section{FT: regulators to tackle shadow banking}\label{sec:di-ft}

\begin{casestudy}[FT, November 2012: ``Regulators to tackle shadow banking'']
A report led by Lord Adair Turner, the UK regulator, calls for regulation of shadow banking; Turner: shadow banking ``is like cholesterol'', there is good and
bad \ts{17}{1}{0:07}. Mehrling, who was working on the subject, thinks shadow banking is in our future and hopes
it will be the good kind \ts{17}{1}{0:28}.
\end{casestudy}

\begin{nfigure}
\begin{taccts}
\tacct[2.4cm]{Traditional (``Jimmy Stewart'') bank}{Loans}{Deposits}\qquad
\tacct[3.2cm]{Shadow bank}{RMBS (tranched)\\Interest rate swap\\Credit default swap}{Repo (RP)\\ABCP held by MMMF}
\end{taccts}
\caption{Traditional and shadow bank, as on Mehrling's slide \ts{17}{1}{1:51}. The traditional bank has a solvency backstop
(FDIC) and a liquidity backstop (Fed); the shadow bank has neither directly, only liquidity puts to traditional banks.}\label{fig:di-banks}
\end{nfigure}

The \textbf{traditional bank}, deposits against loans, is the ``Jimmy Stewart bank'' of the holiday film with a run on his bank. It is backstopped by the FDIC,
which insures the deposits (a solvency backstop), and by the Fed (a liquidity backstop); it is the balance sheet most people have in mind when they say ``bank''
\ts{17}{1}{1:51}. The world is increasingly like the \textbf{shadow bank}: mortgages packaged into residential mortgage-backed
securities, tranched into risky and risk-free tranches, insured by interest rate and credit default swaps, and funded in the money market by repo or asset-backed
commercial paper held by money market mutual funds \ts{17}{1}{2:53}. On both sides of its balance sheet the prices are market prices: the money
market funding rate and the price of risk; hence the coming lectures on swaps \ts{17}{1}{3:34}.

\begin{definition}[New concepts: shadow banking]
\emph{Shadow banking}: money market funding of capital market lending, outside the traditional banking system and its backstops \ts{17}{1}{0:49}
(``shadow'' because it does banking business outside the light of bank regulation).
\end{definition}

\begin{history}[Jimmy Stewart's bank]
In Frank Capra's \emph{It's a Wonderful Life} (1946) James Stewart plays George Bailey, who runs the Bailey Building and Loan in Bedford Falls; during a run he
tells the depositors that their money is not in a vault but in each other's houses.\par
\emph{Reference:} F.~Capra (dir.), \emph{It's a Wonderful Life}, Liberty Films/RKO, 1946.
\end{history}

% ------------------------------------------------------------------
\section{Bagehot's world: two separate markets}\label{sec:di-bagehot}

\paragraph{The money market: indirect finance.} In Bagehot's world banks discounted short-term bills that firms issued to finance goods on their way to final sale;
a portfolio of bills maturing at various dates assured the bank's liquidity \ts{17}{2}{2:06}. New in the course: the bank seen as an
\emph{intermediary}, between a primary lender (the household, for the first time in the course, holding deposits to make payments) and a primary borrower (a
business financing working capital with bills): deposits fund loans. This is \emph{indirect finance} \ts{17}{2}{2:47}.

\paragraph{The capital market: direct finance.} There was also a capital market: government bonds (including perpetual bonds paying a coupon forever), some
corporate bonds for large projects, and equity, held directly by rich people as wealth: \emph{direct finance}, no intermediary
\ts{17}{2}{4:30}. The two markets touch when bonds are issued, since the household pays with deposits that the firm
spends; afterwards the bond market seems separate. Bonds promise small coupons and principal only at maturity, perhaps 30 years later; the staggered-timing
liquidity issue of bills does not arise \ts{17}{2}{5:54}.

\begin{nfigure}
\begin{taccts}
\tacct[2.0cm]{Household (primary lender)}{Deposits\\Bonds, equity}{}\quad
\tacct[2.0cm]{Bagehot bank}{Bills}{Deposits}\quad
\tacct[2.0cm]{Business (primary borrower)}{Working capital\\Fixed capital}{Bills\\Bonds, equity}
\end{taccts}
\caption{Bagehot's world: indirect finance in the money market, direct finance in the capital market \ts{17}{2}{2:47}.}\label{fig:di-bagehot}
\end{nfigure}

\begin{aside}[A cult of liquidity?]
Some economic historians (among them some of Mehrling's professors) argue that Britain's hyper-developed short-term money market, the centre of world trade
finance, came with underdeveloped long-term finance: banks were forbidden by a ``cult of liquidity'' from long-term lending, which may help explain Britain
lagging after the industrial revolution. Mehrling does not vouch for it: ``I'm not an expert on Britain'' \ts{17}{2}{7:15}.
\end{aside}

% ------------------------------------------------------------------
\section{The new world}\label{sec:di-newworld}

The US was a developing country, with an extreme need for long-term finance: railroads, buildings, factories, homes. American banks held corporate bonds and
long-term loans, including ``six-month'' loans that were always rolled over, and few short-term bills \ts{17}{3}{0:03}.
Before there was a central bank, how did they manage? By cash reserves, by bankers' (correspondent) balances, and by interbank borrowing against bonds as
collateral: essentially repo, a paired sale and repurchase \ts{17}{3}{1:26}. A deficit bank
at the clearing pays in cash, in bankers' balances, or borrows against a bond.

\begin{nfigure}
\begin{taccts}
\tacct[2.4cm]{Deficit bank}{Bonds (collateral)\\Loans\\$-$Cash \emph{or} $-$Bankers' balances}{Deposits\\\emph{or} $+$Repo borrowing}\qquad
\tacct[2.4cm]{Surplus bank}{Bonds, loans\\Cash reserves\\$+$Cash \emph{or} $+$Repo lending}{Deposits\\Bankers' balances}
\end{taccts}
\caption{Interbank settlement in the US before the Fed \ts{17}{3}{1:50}.}\label{fig:di-us}
\end{nfigure}

\paragraph{Deficit at the clearing.} A student asks whether, in a developing country with many borrowers, the central bank acted as the surplus bank. There was no
central bank; and ``deficit'' refers to the clearing, not the whole system. Banks lend, as always, by creating deposits; when the borrower spends the deposit
in New York, the bank must pay the New York bank, with reserves, with bankers' balances at its correspondent, or by borrowing against a bond
\ts{17}{3}{3:55}.

\paragraph{The origin of rating agencies.} In Mehrling's account rating agencies arose in this period: banks paying each other did not want to analyse every bond
offered as collateral, so they bought Moody's book of ratings, which said a bond was unlikely to default in the next year, good enough for short-term collateral;
the rating entered the terms of the repo \ts{17}{3}{5:38}. Unsecured lending through bankers'
balances needs a relationship; secured lending against a rated railroad bond can be done with a New York bank that has never heard of you \ts{17}{3}{7:20}.
Capital market and money market are completely intertwined: ``19th-century shadow banking'' \ts{17}{3}{7:41}.

\begin{history}[Moody's, 1909]
John Moody published \emph{Moody's Analyses of Railroad Investments} in 1909, assigning letter ratings to railroad bonds; Poor's (1916), Standard Statistics
(1922) and Fitch (1924) followed. The link to interbank collateral is Mehrling's interpretation.\par
\emph{References:} R.~Sylla, ``An historical primer on the business of credit rating'', in R.~M.~Levich et al.\ (eds), \emph{Ratings, Rating Agencies and the
Global Financial System}, Kluwer, 2002.
\end{history}

\paragraph{The Fed and the shadow system.} After the 1907 crisis the reformers tried to replace this system with discounting of short-term bills at a new Fed, but
not enough bills were issued: the US needed long-term finance \ts{17}{3}{7:41}. Mehrling's argument: the Fed's mistake around 1929
was to let the shadow banking system of its day collapse, refusing to support mortgage credit and bonds as improper for banks, and so it brought down the US
banking system and the world's \ts{17}{3}{8:22}. In 2007 we were fortunate that the Fed did backstop the shadow banking system,
much of it in Europe, through the liquidity puts linking shadow banks to Citibank and other traditional banks \ts{17}{3}{9:03}. The
founders of the Fed ``pasted a kind of central bank onto'' a system it did not fit; today's question is how to have a central bank that supports a system like
this \ts{17}{3}{10:04}. The notes put it sharply: the banking collapse of the 1930s can be read as a run on the shadow banking system of the time, private funding liquidity through repo on
non-Treasury securities; the Fed refused to support it, it collapsed, and the government had to take on long-term lending itself, most famously through the Reconstruction Finance Corporation:
``this mistake was not, thankfully, repeated in 2007--8!'' The transmission from funding liquidity to market liquidity and asset prices is, in the US, ``not recent, but in fact ancient''
\mn{17}{2}\mn{17}{3}.

% ------------------------------------------------------------------
\section{Funding liquidity versus market liquidity}\label{sec:di-shiftability}

The British concept of liquidity was \emph{funding liquidity}: guaranteed cash inflows. In the new world a bank's liquidity was having bonds pledgeable or
shiftable to someone else: \emph{shiftability}, which the course calls market liquidity, the ability to buy or sell an asset quickly without moving the price much
\ts{17}{4}{0:04}.

\begin{definition}[New concepts: shiftability]
\emph{Shiftability}: the liquidity of an asset that can be shifted to another holder (sold or pledged) without loss, as opposed to liquidity from self-liquidating
short-term assets. Term of Harold Moulton, 1918 \ts{17}{4}{0:45}.
\end{definition}

Moulton argued that shiftability was the right concept for the US, and that the Fed should backstop it. The Fed did, not for all assets but for Treasury bills:
that is where open market operations came from, supporting market liquidity in the capital market \ts{17}{4}{1:06}. This course
uses both kinds of liquidity and their interlink, dealers who create market liquidity must fund their inventories, an understanding not yet there in 1918
\ts{17}{4}{1:46}.

\begin{history}[Harold G.~Moulton]
Harold Glenn Moulton (1883--1965) published ``Commercial Banking and Capital Formation'' in four parts in the \emph{Journal of Political Economy} in 1918; he
was the first president of the Brookings Institution (1927--1952).\par
\emph{References:} H.~G.~Moulton, \emph{JPE} vol.~26, 1918 (four parts); Brookings Institution, ``History''.
\end{history}

% ------------------------------------------------------------------
\section{Schumpeter on banking}\label{sec:di-schumpeter}

Joseph Schumpeter, sympathetic to the American scene, emphasized all his life the importance of banks willing to do long-term development finance: money market
funding of capital market lending \ts{17}{5}{0:04}. Banks create new money for entrepreneurs, who need not find
someone who has saved: ``created from nothing'', the alchemy of banking. To say so risks sounding like a monetary crank, and Schumpeter was not one
\ts{17}{5}{1:08}. Expanding the balance sheet is not ``bread from heaven'': something must discipline it (next section) \ts{17}{5}{1:50}.

Schumpeter admired H.~D.~Macleod, who stressed this kind of banking in the development of Scotland. The Lewis model of development, with surplus labour in
subsistence agriculture, also gives banking a role in mobilizing it towards cities and factories \ts{17}{5}{2:10}.
Development economists are usually ``real side'' people; Mehrling would give half of a development course to finance, ``but hey, I would say that''
\ts{17}{5}{3:32}.

\begin{intuition}[There is such a thing as a free lunch]
The notes warn that most self-styled monetary alchemists are cranks: ``the ability to create money from nothing is not the same as the ability to create bread from nothing'', and economics has
come close to making ``there ain't no such thing as a free lunch'' a creed. ``But there is.'' Adam Smith (1776) already urged paper money to economize on gold, trading a sterile commodity for
productive capital; the main source of the free lunch, though, is that banking can \emph{mobilize unused resources}: if entrepreneurs spend new deposits on things that would otherwise go unsold,
they raise activity, not prices \mn{17}{3}.
\end{intuition}

\begin{coursenote}
In class Schumpeter's ``thesis'' is dated 1912. The book meant is \emph{Theorie der wirtschaftlichen Entwicklung} (published late 1911, dated 1912; English
\emph{The Theory of Economic Development}, 1934), not his doctoral thesis. Henry Dunning Macleod (1821--1902) was a Scottish economist, author of \emph{The Theory
and Practice of Banking} (1855--56). The Lewis model is W.~A.~Lewis, ``Economic Development with Unlimited Supplies of Labour'', \emph{Manchester School}, 1954,
which discusses capital formation financed by credit creation \ts{17}{5}{0:04}.
\end{coursenote}

% ------------------------------------------------------------------
\section{Payment versus funding}\label{sec:di-payment}

A business wants to buy a machine from ``society''. The bank lends by a swap of IOUs, $+$loan against $+$deposit, and the business pays with the deposit
\ts{17}{6}{0:05}. The question: what \emph{funds} the loan? If society is happy to add the
deposit to its portfolio, the deposit funds the machine, even though it is a demand deposit and the machine a loom with a ten-year life
\ts{17}{6}{1:53}. If society does not want to hold it, the business must fund the machine otherwise, by issuing a bond that
society buys with the deposit, and repaying the loan \ts{17}{6}{2:55}.

\begin{nfigure}
\begin{taccts}
\tacct[2.0cm]{Bank}{$+$Loan\\(later $-$Loan)}{$+$Deposit\\(later $-$Deposit)}\quad
\tacct[2.0cm]{Business}{$+$Machine}{$+$Loan\\(later $-$Loan, $+$Bond)}\quad
\tacct[2.0cm]{Society}{$-$Machine, $+$Deposit\\(later $-$Deposit, $+$Bond)}{}
\end{taccts}
\caption{Payment versus funding: bank credit as a bridge loan, the bond as take-out financing \ts{17}{6}{1:10}.}\label{fig:di-payment}
\end{nfigure}

The money market intermediation is a bridge loan: it gets the ball rolling, using the elasticity of the payment system; once the machine is running, the capital
market provides permanent, take-out financing. First intermediation, as in Bagehot but long; then direct finance, bond for bond, the intermediary drops out
\ts{17}{6}{4:17}.

% ------------------------------------------------------------------
\section{Gurley and Shaw: intermediaries in the capital market}\label{sec:di-gurleyshaw}

\begin{reading}[J.~G.~Gurley and E.~S.~Shaw, Money in a Theory of Finance (1960)]
Gurley and Shaw (Brookings, 1960) introduced the terms \emph{direct finance} (the primary lender lends directly to the primary borrower) and \emph{indirect
finance} (an intermediary in between). They stressed intermediation not just in the money market but in the capital market, for long-term finance, following
Moulton; and not only by banks but by the other intermediaries that had grown since his day \ts{17}{7}{0:04}.
\end{reading}

\begin{definition}[New concepts: direct and indirect finance, financial intermediary]
\emph{Direct finance}: the primary borrower's securities are held by the primary lender. \emph{Indirect finance}: a \emph{financial intermediary} (from Latin
\emph{inter}, between, and \emph{medius}, middle) issues its own liabilities to lenders and holds the borrowers' securities \ts{17}{7}{0:31}.
\end{definition}

% ------------------------------------------------------------------
\section{Financial evolution}\label{sec:di-finevol}

The two main capital market intermediaries in the US are pension funds and insurance companies \ts{17}{8}{0:05}:
\begin{itemize}
  \item a \textbf{pension fund}, typically a defined benefit plan in 1960, promises workers a fraction of final income for life, a long-term, annuity-like promise,
        and hedges it with long-term assets, especially equities \ts{17}{8}{0:46};
  \item an \textbf{insurance company} has contingent liabilities (fire, life policies) and holds bonds \ts{17}{8}{1:30}.
\end{itemize}
Both stand between households and businesses, like Bagehot's bank, though their liabilities are not money; they became bigger than banks in funding American
development \ts{17}{8}{2:11}.

\begin{intuition}[Gurley and Shaw's vision: intermediaries face both ways]
Households want to hold liquid assets, money, insurance, pensions; businesses want to borrow long. How can the capital development of the nation be financed?
An intermediary faces both ways: to households it says ``I'll give you what you want'', deposits, policies, pensions; to businesses, ten-year loans, or it buys
their bonds and equity. Without it, the mismatch between what households want and what businesses issue could block development
\ts{17}{8}{3:33}.
\end{intuition}

Their liquidity: insurance has an actuarial quality (not every house burns at once), and in a hurricane the insurer sells bonds; pension funds turn equities into
bonds (annuities) when a worker retires. Both depend on market liquidity, shiftability, as Moulton said for banks; Moulton became president of Brookings, which
may be why Brookings hired Gurley and Shaw \ts{17}{8}{5:17}.

\paragraph{The flaw: risk does not go away.} The risk of the assets does not disappear because you promise something else: it falls on the liabilities, or on
whoever guarantees them, typically the government. With risky assets, guaranteed benefits may not be paid; the risk is borne by the sponsor (often a corporation),
by the government through the Pension Benefit Guaranty Corporation, or by the beneficiaries when promises are not kept. Finance and arbitrage will make you
appreciate it, through mispricing \ts{17}{8}{7:20}.

\paragraph{From indirect to direct finance.} Since 1960 indirect finance has been hollowed out in favour of mutual funds: stock funds hold equities and issue shares
in the portfolio (Fidelity's Magellan), bond funds hold bonds (PIMCO): the risk passes right through to the holder, ``this is not investment advice''
\ts{17}{8}{9:07}. Pensions moved from defined benefit to defined contribution plans such as the
401(k), where you see your risk \ts{17}{8}{10:51}. In insurance, term policies replaced whole-life policies with their savings component: both trends reflect mutual funds competing
with traditional indirect finance. A mutual fund looks like an intermediary, but it mainly pools risk through diversification, with little transformation: its shares have the risk of the
underlying pool; the liquidity of buying back at NAV requires cash or credit lines, paid for by the shareholders. Bond funds replaced the savings in whole-life policies and bank time deposits,
stock funds replaced defined-benefit plans \mn{17}{6}\mn{17}{7}. Gurley and Shaw ``overdid it'', the notes say: modern finance stresses that intermediation eliminates no risk, only transfers
it, sometimes opaquely, so capital markets equilibrate by price, and attempts to avoid this create arbitrage opportunities that have transformed the system \mn{17}{6}.

\begin{keyformulas}
\textbf{Quantities versus prices} \ts{17}{8}{11:32}. Indirect finance solves the mismatch between what lenders want and what
borrowers want with \emph{quantities}, on its own balance sheet. Modern finance solves it with \emph{prices}: if you don't want to hold bonds, bonds fall in
price until you do. Price does much more work than in 1960, which is why everyone watches the stock market.
\end{keyformulas}

The same evolution, from intermediation to direct finance, has happened in banking and the money market: shadow banking is its natural next step \ts{17}{8}{12:33}.

% ------------------------------------------------------------------
\section{Banking evolution}\label{sec:di-bankevol}

Banks as intermediaries are a new conception in the course, which has treated banking as a payment system and as market making, not as funding
\ts{17}{9}{0:06}. Why can banks fund things? Because people want deposits not just to pay but as an asset, a liquidity reserve in their wealth
\ts{17}{9}{0:26}:
\begin{itemize}
  \item deposits are \textbf{permanent} funding for the system as a whole: money demand fluctuates but does not go to zero, so reserves need cover only the
        fluctuations; a banker who sees deposits never fall below 80\% of the balance sheet can hold 80\% in earning loans \ts{17}{9}{2:10};
  \item deposits are \textbf{cheap} funding, close substitutes for cash that pays no interest; even when banks compete by paying interest, it is less than on bonds
        \ts{17}{9}{1:30}.
\end{itemize}

\paragraph{Liquidity at the margin.} In the notes: because households and firms want to \emph{hold} means of payment, not only make payments, the banking system can keep a permanent
short position in cash and use it to fund long positions in non-cash assets. Liquid liabilities need liquid assets only at the margin: cash reserves plus the ability to replenish them (Fed funds,
the discount window, ``secondary reserve'' assets that can be sold); the rest of the deposits can go into less liquid assets. Once the focus is intermediation, solvency risk, interest rate risk and
credit risk, comes to the fore \mn{17}{7}\mn{17}{8}.

\paragraph{Sharing out the spoils.} Everyone wants this cheap funding \ts{17}{9}{3:52}: the government through the central bank (Treasury bills
against cash), business through commercial banks (industrial loans against deposits), households through savings and loans, the Jimmy Stewart banks (mortgages
against ``deposits'', legally shares but used as money) \ts{17}{9}{4:33}.

\begin{nfigure}
\begin{taccts}
\tacct[1.9cm]{Central bank}{Treasury bills}{Cash}\quad
\tacct[1.9cm]{Commercial bank}{Business loans}{Deposits}\quad
\tacct[1.9cm]{Savings and loan}{Household mortgages}{``Deposits'' (shares)}
\end{taccts}
\caption{Cheap permanent funding shared out among government, business and households \ts{17}{9}{4:33}.}\label{fig:di-spoils}
\end{nfigure}

\paragraph{The risk passes through, to the taxpayer.} The banks face both ways, but the risk does not vanish: in the Jimmy Stewart bank the FDIC guarantees
solvency and the Fed liquidity, so the taxpayer bears it. Hence regulation (no risky loans, no interest on deposits), and hence a regulatory arbitrage to create
non-banks doing the same thing: that is where shadow banking came from \ts{17}{9}{6:17}. In
the Jimmy Stewart bank the risk stays in the community: ``your money is in your houses'' \ts{17}{9}{7:40}.

\paragraph{Shadow banking strips the risks.} In shadow banking the loans become traded securities whose price is the price of all their risks: a mortgage has
duration risk (30 years) and default risk (household income). The interest rate swap and the credit default swap strip them off, leaving a synthetic Treasury
bill funded in the money market: risk is not eliminated but hedged, sold off separately, rather than borne by a government that pretends it is not there, with the
moral hazard that implies \ts{17}{9}{8:22}. As pension funds evolved into stock funds and insurers into bond funds,
banks evolved into money market mutual funds, many outside the US, holding dollars for sovereign wealth funds and central banks: ``a global dollar system funding
subprime mortgages in Cleveland'' \ts{17}{9}{9:44}. From indirect finance (Gurley and Shaw, 1960) to direct finance (2012),
with all risks priced in markets and money and capital markets completely integrated \ts{17}{9}{10:45}.

% ------------------------------------------------------------------
\section{Preview: central banking and shadow banking}\label{sec:di-preview}

\paragraph{The central bank's exposure.} A student asks how the central bank becomes exposed to shadow banks \ts{17}{10}{0:10}. To the traditional
bank the Fed gives a liquidity backstop only; if it is insolvent the FDIC shuts it down and pays depositors from accumulated premia, and beyond that the
government backstops solvency \ts{17}{10}{0:52}. Nobody directly backstops the shadow bank, but there is the liquidity put: when the quality of the
collateral came into question, money funds refused to roll their funding (``you told me these were just like Treasury bills with a yield advantage''), and the
shadow banks drew on their parents' promises, say Citibank's, to roll their funding; so they came back onto traditional banks' balance sheets
\ts{17}{10}{1:33}. Regulators had endorsed the system to get risk out of banks: it removed most solvency risk but not liquidity risk
\ts{17}{10}{2:35}.

\begin{definition}[New concepts: liquidity put]
A commitment by a bank to provide funding to an off-balance-sheet vehicle if it cannot roll its money market funding: in effect an option to ``put'' the funding
problem back to the bank \ts{17}{10}{1:33}.
\end{definition}

\paragraph{Where is solvency?} Most banking courses are about solvency; this course stresses liquidity. In the shadow bank the solvency risk sits with the holders of
the swaps and of the mezzanine and junior tranches. For AAA credit default swaps it was AIG, an insurance company with a lot of capital and bonds, the ``FDIC
equivalent'' of the system, whose capital was sucked out very quickly by the counterparties, Soci\'et\'e G\'en\'erale and Goldman Sachs and their clients, which were
acting as market makers \ts{17}{10}{2:55}.

\begin{intuition}[The fantasy of modern finance]
The shadow banking system had no explicit lender of last resort, on the belief that removing solvency risk removes liquidity risk: a risk-free asset is like a
Treasury bill, hence liquid by design. False: someone must still buy it. Market liquidity is created by dealers quoting bid--ask spreads, and dealers must fund
their inventories; if they cannot, they stop making markets, prices go away, nobody knows what the collateral is worth, and the ability to fund the portfolio goes
away. That is what happened in 2007--09 \ts{17}{10}{5:03}.
\end{intuition}

A traditional bank has no market prices on either side: deposits were forbidden to pay interest, and loans are illiquid, specific to each borrower. A shadow bank
lives on two prices, the market price of the RMBS and the world dollar funding rate \ts{17}{10}{6:26}. The system had moved off
the Fed's balance sheet, ``out of sight, out of mind'', Europe's problem; wrong: a global funding system came under threat almost immediately, and the expansion of
the Fed's balance sheet shown in the first lecture bailed it out \ts{17}{10}{7:49}.

% ------------------------------------------------------------------
\flushside
\section*{Key formulas and ideas}
\begin{keyformulas}
\begin{tabularx}{\linewidth}{@{}l L@{}}
Shadow banking & money market funding of capital market lending; prices on both sides\\
Bagehot's world & money market (bills, indirect) and capital market (bonds, equity, direct) separate\\
New world & long-term bank assets, interbank repo against rated bonds: capital and money markets intertwined\\
Liquidity & funding liquidity (British) versus shiftability $=$ market liquidity (Moulton 1918)\\
Payment vs funding & bank credit as bridge loan; bonds as take-out finance\\
Gurley--Shaw & intermediaries face both ways; but risk passes through to someone\\
Evolution & intermediation $\to$ mutual funds; banks $\to$ MMMFs and shadow banks; quantities $\to$ prices\\
Backstops & traditional: FDIC (solvency), Fed (liquidity); shadow: liquidity puts, AIG as solvency\\
\end{tabularx}
\end{keyformulas}

\section*{Exercises}

\begin{exercise}[Bridge and take-out]
A firm borrows 100 from a bank to buy a machine and, a year later, issues a 10-year bond of 100 bought by a household. Write the T-accounts of bank, firm and
household at each step, and say which step is indirect and which direct finance.
\end{exercise}

\begin{exercise}[Permanent funding]
A bank's deposits have fluctuated between 800 and 1000 over ten years. How much of its balance sheet can it hold in long-term loans, and how much in reserves?
What changes if a money market fund offers a close substitute for deposits?
\end{exercise}

\begin{exercise}[Who bears the risk?]
For a defined benefit pension fund, a defined contribution 401(k), a Jimmy Stewart bank and a shadow bank holding an RMBS hedged with swaps, say who bears the
default risk and who the liquidity risk.
\end{exercise}

\begin{exercise}[Market liquidity needs funding]
Using Treynor's model, explain why a dealer in AAA mortgage securities stops quoting when its repo funding is cut, and what happens to the value of the same
securities as repo collateral.
\end{exercise}
